\chapter*{Course Overview}
\addcontentsline{toc}{chapter}{Course Overview}

This 90-day curriculum is designed for students preparing to enter a Computer Science PhD program in AI, large language models, vision-language models, reinforcement learning, world models, and multimodal reasoning. The curriculum follows the historical development of modern foundation models while emphasizing engineering practice, experimental design, pretraining experience, reinforcement learning foundations, world model research, scalable inference, and research-oriented implementation.

Each chapter is written as a self-contained unit that defines what to study, why it matters, how to implement it, where things go wrong, and what to write about. Students should expect to read papers, run experiments, maintain reproducible logs, and produce short research memos throughout the program.

\chapter*{Learning Goals}
\addcontentsline{toc}{chapter}{Learning Goals}

By the end of this curriculum, the student should be able to:

\begin{itemize}
    \item Understand the architectural evolution from Transformer to modern LLMs and VLMs.
    \item Configure, train, and diagnose a decoder-only Transformer using modern frameworks.
    \item Understand tokenizer internals and their impact on efficiency, multilinguality, code, math, and reasoning.
    \item Build and debug an end-to-end causal language model pretraining pipeline.
    \item Design a pretraining data pipeline with filtering, deduplication, decontamination, quality scoring, and mixture control.
    \item Pretrain a 300M-scale language model from scratch under a realistic token budget.
    \item Understand the practical role of distributed training, memory sharding, parallelism, and throughput measurement.
    \item Fine-tune and preference-optimize 1B--3B models.
    \item Understand instruction tuning, DPO, RLHF, RLVR, GRPO, and verifiable reward.
    \item Configure and diagnose core reinforcement learning algorithms from tabular methods to PPO.
    \item Understand exploration, intrinsic reward, curiosity, RND, and information-gain-style methods.
    \item Configure, train, and diagnose world models and model-based agents.
    \item Understand Dreamer-style, JEPA-style, transformer-based, diffusion-based, and MPC-style world models.
    \item Build retrieval systems and distinguish naive RAG from agentic retrieval.
    \item Deploy models for scalable evaluation using inference engines such as vLLM or SGLang.
    \item Design evaluations with clear metrics, baselines, statistical uncertainty, contamination checks, and human-evaluation protocols.
    \item Understand long-context methods and their limitations.
    \item Analyze VLM hallucination, grounding failures, and intervention methods.
    \item Apply mechanistic intervention and representation engineering methods where appropriate.
    \item Produce a capstone research report with baselines, ablations, error analysis, and limitations.
\end{itemize}

\chapter*{Course Structure}
\addcontentsline{toc}{chapter}{Course Structure}

\begin{itemize}
    \item \textbf{Part I: Transformer Foundations and Tokenization} \hfill Days 1--12
    \item \textbf{Part II: Data, Pretraining, Distributed Training, and Scaling} \hfill Days 13--30
    \item \textbf{Part III: Post-Training} \hfill Chapters 11--17
    \item \textbf{Part IV: World Models and Model-Based Agents} \hfill Chapters 18--21
    \item \textbf{Part V: Evaluation, VLMs, Serving, Long Context, and Intervention} \hfill Chapters 22--27
    \item \textbf{Part VI: Research Capstone} \hfill Chapter 28
\end{itemize}

\chapter*{Research Workflow}
\addcontentsline{toc}{chapter}{Research Workflow}

This curriculum assumes you will use AI coding assistants---Claude, Copilot, ChatGPT, or similar tools---as your primary implementation partner. This is not a concession; it reflects how productive researchers actually work today. You will not be evaluated on whether you can write a training loop from memory. You will be evaluated on whether you can:

\begin{enumerate}
    \item \textbf{Specify} what needs to be built---the right architecture, the right loss, the right data pipeline for the research question at hand.
    \item \textbf{Verify} that what was built is correct---by designing sanity checks, tracing gradient flow, inspecting intermediate outputs, and catching silent bugs before they corrupt your results.
    \item \textbf{Diagnose} why something does not work as expected---by reading loss curves, gradient norms, attention patterns, and generated samples, and forming testable hypotheses about the root cause.
    \item \textbf{Design experiments} that answer a research question---with proper baselines, controlled ablations, appropriate metrics, and honest reporting of negative results.
    \item \textbf{Interpret results} with scientific rigor---distinguishing real effects from noise, identifying confounders, and knowing when a conclusion is (or is not) supported by the evidence.
\end{enumerate}

Every lab in this curriculum follows a four-stage structure that mirrors a real research workflow:

\begin{itemize}
    \item \textbf{Setup.} We provide a working-but-imperfect baseline or a clear specification for one. You start from running code, not a blank file.
    \item \textbf{Build.} You extend the baseline---adding modules, metrics, or configurations. Use an AI assistant freely; the point is getting to a runnable experiment, not hand-typing boilerplate.
    \item \textbf{Verify.} Before running any experiments, you design and execute sanity checks to confirm the implementation is correct. This is the most underrated research skill.
    \item \textbf{Investigate.} You run controlled comparisons, record results, and write a short analysis. The deliverable is an experiment report, not a code submission.
\end{itemize}

\chapter*{Execution Policy}
\addcontentsline{toc}{chapter}{Execution Policy}

The curriculum has three workload tiers. The \textbf{Core} tier is required for every student and should fit the 90-day schedule for a full-time learner. The \textbf{Optional} tier deepens skill when compute and time permit. The \textbf{Stretch} tier is intended for students with strong systems background, extra GPU access, or a capstone directly tied to the topic.

\begin{itemize}
    \item \textbf{Core:} read the key papers, complete the lab (all four stages), and write the experiment report.
    \item \textbf{Optional:} run additional ablations, reproduce a public baseline, or extend the investigation.
    \item \textbf{Stretch:} scale the experiment, compare multiple models, add distributed training, or turn the result into capstone-grade evidence.
\end{itemize}

\chapter*{Weekly Milestones}
\addcontentsline{toc}{chapter}{Weekly Milestones}

\begin{itemize}
    \item \textbf{Weeks 1--2:} working mini-GPT implementation, tokenizer comparison, and reproducible training loop.
    \item \textbf{Weeks 3--4:} data pipeline, pretraining dry run, scaling-law analysis, and distributed-training notes.
    \item \textbf{Weeks 5--6:} SFT, PEFT, DPO or toy RLHF, and a first aligned-assistant evaluation report.
    \item \textbf{Weeks 7--8:} evaluation methodology, one focused VLM/RAG/long-context experiment, serving setup only if needed, and a VLM benchmark report.
    \item \textbf{Weeks 9--10:} RL foundations, PPO baseline, intrinsic-reward experiment, and RL failure analysis.
    \item \textbf{Weeks 11--12:} latent dynamics model, model-based RL comparison, world-model report, capstone proposal, and capstone implementation freeze.
    \item \textbf{Week 13:} capstone ablations, final report, reproducibility package, and presentation.
\end{itemize}

\chapter*{Chapter Format}
\addcontentsline{toc}{chapter}{Chapter Format}

Each chapter follows the same structure:

\begin{itemize}
    \item Conceptual core and historical motivation
    \item Mathematical intuition (key equations with emphasis on \emph{what they mean}, not derivation for its own sake)
    \item Trick table (5 or fewer operationally critical techniques, with expected effects and failure signals)
    \item Failure modes and common misconceptions
    \item Lab: Setup / Build / Verify / Investigate
    \item Mini experiment report (replaces the traditional ``research memo'')
    \item Key readings with guided annotations
\end{itemize}

Interspersed throughout the conceptual sections are \emph{intuition checkpoints}---short questions designed to test whether you understand what an equation or mechanism is \emph{doing}, not whether you can re-derive it. If you cannot answer a checkpoint from understanding alone, re-read the preceding section before moving on.

Every experiment report is \core{} because the course is meant to build research judgment, not only implementation fluency.

\chapter*{Key Reading Map}
\addcontentsline{toc}{chapter}{Key Reading Map}

\begin{itemize}
    \item \textbf{Chapters 1--5:} Bengio et al.\ on neural language models; Mikolov et al.\ on Word2Vec; Bahdanau et al.\ on attention; Vaswani et al.\ on Transformers; Sennrich et al.\ on BPE; Kudo and Richardson on SentencePiece; Radford et al.\ on GPT.
    \item \textbf{Chapters 6--7:} Devlin et al.\ on BERT; Radford et al.\ on GPT-2; Brown et al.\ on GPT-3 and in-context learning; Min et al.\ and Xie et al.\ on in-context learning analysis.
    \item \textbf{Chapters 8--10:} Penedo et al.\ on RefinedWeb; Together AI on RedPajama; Dodge et al.\ on C4 documentation; Kaplan et al.\ and Hoffmann et al.\ on scaling laws; Rajbhandari et al.\ on ZeRO; PyTorch FSDP documentation; Narayanan et al.\ and Shoeybi et al.\ on large-model parallelism.
    \item \textbf{Chapters 11--14:} Wei et al.\ on instruction tuning; Ouyang et al.\ on InstructGPT; Hu et al.\ on LoRA; Dettmers et al.\ on QLoRA; Rafailov et al.\ on DPO; Schulman et al.\ on PPO; DeepSeek-AI and related reasoning-RL reports for RLVR and GRPO-style training.
    \item \textbf{Chapters 15--20:} Liang et al.\ on HELM; Srivastava et al.\ on BIG-bench; Zheng et al.\ on MT-Bench and Chatbot Arena; Dosovitskiy et al.\ on ViT; Radford et al.\ on CLIP; Li et al.\ on BLIP-2; Liu et al.\ on LLaVA; Kwon et al.\ on vLLM; Su et al.\ on RoPE; Liu et al.\ on Lost in the Middle; Lewis et al.\ on RAG; POPE, MME, and CHAIR papers for VLM hallucination.
    \item \textbf{Chapters 21--24:} Sutton and Barto for RL foundations; Mnih et al.\ on DQN; Schulman et al.\ on policy gradients, GAE, and PPO; Pathak et al.\ on curiosity; Burda et al.\ on RND; Ecoffet et al.\ on Go-Explore.
    \item \textbf{Chapters 25--29:} Sutton on Dyna; Janner et al.\ on MBPO; Hafner et al.\ on Dreamer; Micheli et al.\ on IRIS; Alonso et al.\ on DIAMOND; Hansen et al.\ on TD-MPC2; LeCun on JEPA; recent VLM-agent papers should be selected by the instructor based on the capstone direction; Heilmeier's catechism; NeurIPS reproducibility checklist.
\end{itemize}
